% sections/07_conclusion.tex
% §7 Conclusion — RW (Võ Thanh Phong, SE194101)
% Nguồn: RQ1_results_FINAL.md + full_metrics_output_FINAL.csv + proposal.md §4

\section{Conclusion}
\label{sec:conclusion}

We investigated whether Qwen2.5-7B-Instruct with 3-shot prompting
can match the quality thresholds established by GPT-4o 3-shot
(CTQRS$\geq$75\%, SBERT$\geq$0.82, ROUGE-1$\geq$0.47)
on 262 Bugzilla bug reports~\cite{acharya2025bugquality}.
Our results are mixed. Qwen2.5-7B few-shot exceeds the
structural quality threshold (median CTQRS=88.89\%,
$p<0.001$, $\delta=0.908$, large) and the lexical overlap
threshold (median ROUGE-1=0.491, $p=0.002$, $\delta=0.191$,
small), but falls substantially below the semantic
similarity threshold (median SBERT=0.632, $p=1.000$,
$\delta=-0.969$, large). An ablation study adding a
length constraint to the prompt improved ROUGE-1
(0.491$\rightarrow$0.532) but further reduced SBERT
(0.632$\rightarrow$0.609), confirming that the SBERT
deficit stems from output style mismatch rather than
semantic content errors.

This is the first study to directly evaluate Qwen2.5-7B
few-shot prompting against published GPT-4o few-shot thresholds
on the bug report restructuring task. Our findings suggest
that, for organizations with limited GPU resources that
cannot afford fine-tuning, Qwen2.5-7B 3-shot prompting is
a viable option for producing structurally complete bug reports
(CTQRS) with adequate lexical fidelity (ROUGE-1), but style
alignment with Bugzilla-style ground-truth reports remains
an open challenge that few-shot prompting alone cannot resolve.

Future work could investigate three concrete directions:
(1) applying style-specific constraints or in-context style
exemplars that more precisely match the plain-text, concise
Bugzilla writing style to close the SBERT gap without
fine-tuning;
(2) cross-domain evaluation on Mojira/Minecraft or mobile
app bug trackers to assess generalizability beyond Mozilla
projects;
(3) lightweight parameter-efficient fine-tuning (e.g.,
LoRA~\cite{acharya2025bugquality}) on a small domain-specific
sample to quantify the minimum training investment needed
to close the remaining gap.

Our replication package --- including the dataset split,
3-shot prompt template, inference scripts, metric scorer,
and raw results --- is publicly available at
\url{https://github.com/HoangThong05/RT-SWT-003-nhom3}.
